% 第 7 章　生活中的力（完整样章）
\chapter{生活中的力}

\step{反常识开场}

\begin{mainline}
夏天的雷雨云底离地面大约两千米。一颗雨滴从那样的高度掉下来，如果中途什么都不拦它，只受重力牵着一路加速，那么到达你伞面时的速度应该是
\[
v=\sqrt{2gh}\approx\sqrt{2\times9.8\times2000}\ \text{米每秒}\approx 200\ \text{米每秒}.
\]
两百米每秒是什么概念？手枪子弹出膛的速度大约就是两三百米每秒。如果雨真以这个速度砸下来，每一颗雨滴都是一颗小子弹，下雨天出门等于走进靶场。

可是你昨天还在雨里散步。雨滴打在伞上噼啪作响，打在脸上有点凉，仅此而已。气象观测给出的答案是：一颗普通雨滴落地的速度大约只有七米每秒——和自由落体预言的两百米每秒差了将近三十倍。

那三十倍的速度到哪里去了？是谁在半途“接住”了雨滴？不是云，不是风，而是一种看不见、却每秒钟都在你身上推推搡搡的力：空气阻力。

再把视线收回室内。你用尽力气推书柜，书柜纹丝不动。第 6 章刚刚说过，合力为零才没有加速度——可你明明在推它，你的手明明酸了，那个把你的推力“顶回去”的力又是谁发的？地板看起来什么都没做，它甚至连动都没动一下。

雨滴、书柜，还有每天托着你不让你陷进地里去的椅子、拉着你转弯的车轮、绷紧的琴弦——它们背后是同一张清单。第 6 章给了三条总规则，告诉你“合力决定运动怎么变”；这一章要回答的是那个被推迟了的现实问题：\emph{具体到每一天，力到底有哪几种，每一种的脾气是什么样的？}
\end{mainline}

\step{先猜一猜}

在翻开答案之前，先押三注。写错比写对更有用。

第一题。书柜静止在地上，你用五十牛顿的力水平推它，它不动。这时地板给书柜的静摩擦力有多大？（a）零，因为书柜没动；（b）小于五十牛顿；（c）恰好五十牛顿；（d）大于五十牛顿，否则书柜会往后滑。

第二题。一根橡皮筋挂一枚硬币，被拉长了一厘米。挂三枚同样的硬币，大约会被拉长多少？（a）还是一厘米，橡皮筋又不认识硬币；（b）两厘米左右，越拉越费劲；（c）三厘米左右，拉多长和挂多重成正比；（d）九厘米，效果会叠加得越来越快。

第三题。下雨天，大颗的雨滴和小颗的雨滴从同一朵云里落下来，哪一种落得更快？（a）一样快，第 5 章说过轻重物体下落一样快；（b）大的快；（c）小的快；（d）看心情。

第四题。汽车急刹车时，有经验的司机会“点刹”——不让车轮完全抱死打滑。这样做最主要的好处是什么？（a）减少轮胎磨损；（b）抱死之后制动力反而变小，而且失去转向；（c）让刹车片冷却；（d）纯属驾驶习惯，没有物理依据。

把三个答案写在纸上。本章结尾我们逐一对账——尤其是第一题，它的答案会让很多人第一次意识到：有一种力会“看情况办事”。

\step{动手实验}

\begin{experiment}[橡皮筋的脾气与斜坡上的书]
\textbf{材料}：一根粗细均匀的橡皮筋、十几枚同样的硬币、一把尺子、一支笔、一张纸、一本厚书、一块尽量平整的木板（或结实的硬纸板）。

\textbf{第一部分：橡皮筋有多守规矩。}把橡皮筋的一端挂在桌沿的固定点上（或者用左手捏住悬空），下端挂一枚硬币，用尺子量出橡皮筋的长度，记在纸上。然后依次挂两枚、三枚、四枚……每加一枚记一次长度，算出每次的伸长量。在纸上画一张图：横轴是硬币枚数，纵轴是伸长量。你会发现什么？点几乎落在一条过原点的直线上——至少在硬币不多的时候是这样。

\textbf{第二部分：把它拉过头。}继续加硬币，或者直接用手把橡皮筋拉到原来的两三倍长，然后取掉所有硬币。量一量它现在的“原长”。它比实验开始时\emph{变长了}，而且再也缩不回去。你在图上把这一段也画出来：直线拐了弯，而且不回头。橡皮筋的“守规矩”是有额度的。

\textbf{第三部分：斜坡上的书。}把厚书放在木板一端，慢慢抬高木板的这一端，让斜坡越来越陡。记下书\emph{刚开始滑动}时的大致角度，用书脊比着在纸上画一条线。然后让书重新开始滑——你会发现，一旦滑起来，木板角度稍微放小一点，书仍然继续滑。也就是说，“刚要滑”和“已经在滑”是两种不同的状态，摩擦在这两种状态下开出的价码不一样。

\textbf{第四部分：纸团与平纸。}取两张完全相同的纸，一张用力揉成紧实的纸团，一张保持平展。双手举到同一高度同时松手。纸团几乎直直落地，平纸却晃晃悠悠、明显慢半拍。两张纸一样重，重力完全相同，差别只在迎风面积——这是空气阻力在家里最便宜的演示。如果教室里大家都做，还可以试试把平纸换成对折再对折的小纸片，看它能快多少。

\textbf{想一想}：第二部分里橡皮筋为什么回不了头？第三部分里，“刚要滑”时的摩擦和“已经在滑”时的摩擦，哪个更大？第四部分里，如果把这个实验搬到没有空气的月球上做，结果会怎样？
\end{experiment}

\step{旧模型遇到麻烦}

第 6 章留下的工具其实非常省：三条规则，外加一句“合力决定动量怎么变”。可是你拿着它走进厨房，立刻会撞上三堵墙。

第一堵墙：\emph{规则没给清单}。“合力”两个字好写，可具体到一只书柜、一颗雨滴，力到底从哪来？重力算一个，然后呢？椅子托着你，这个“托”多大？绳子拉着你，这个“拉”多大？没有一张清单，牛顿第二定律就是一把没有子弹的枪。

第二堵墙：\emph{有的力好像会读心术}。你轻轻推书柜，书柜不动；你加倍用力推，书柜还是不动。按旧直觉，力是施力者“发”出来的，发多少是多少。可地板偏偏不是：你发五十牛，它顶五十牛；你发八十牛，它顶八十牛。它不看你用了多大力，只看\emph{需要}多大力才能不让书柜动。一个会“按需要开价”的力，跟“力是发出来的固定东西”这个老观念正面冲突。这就是本章要修理的核心直觉之一。

第三堵墙：\emph{有的力越跑越强}。骑车的人都知道：慢骑轻松，快骑费劲，而且费劲的程度长得比速度还快。亚里士多德当年说“力维持速度”，我们已经把他请下了台；可空气阻力偏偏给人一种“他说对了一半”的错觉——要维持高速，你确实得不停地蹬。区别在于：你蹬车对抗的不是“运动本身”，而是一个\emph{随速度长大的阻力}。旧模型里没有这种力的位置，清单上必须给它补一行，而且它的“价格表”和前面几种完全不同。

这三堵墙指向同一件事：我们需要的不是更多的总规则，而是一本\emph{力的性格档案}——每种力从哪里来、朝哪个方向、大小由什么决定、什么时候耍赖不认账。

\step{发明新概念}

像清点工具箱一样，我们把日常生活里最常见的力一个个请上台。注意每位的“来源、方向、大小由谁定”这三项档案。

\textbf{重力。}老朋友了：地球对物体的吸引。方向竖直向下，大小在地面附近近似是质量乘以一个常数 \(g\approx9.8\) 牛每千克。它是清单里唯一不挑对象的成员——只要你在地球附近，它就一直在场。第 8 章会告诉你它和月球的运动原来是同一件事。

\textbf{支持力。}椅子托你、桌面托杯子、地板托书柜的那个“托”。来源是\emph{接触面的微小形变}：你坐上去，椅面被压得凹下去一点点，凹下去的材料像被压紧的无数根微型弹簧一样往回顶。方向总是垂直于接触面，所以它在更正式的场合叫“法向力”，大小写时常记作 \(N\)。它的大小不由它自己决定，而由“需要多少才能不让你穿过去”决定——坐个大象上来，它就顶到大象的体重，直到椅子散架为止。

\textbf{拉力（张力）。}绳子、缆线、琴弦的力。来源同样是形变：绳子被拉长一丁点（肉眼往往看不出），内部的分子被扯开，于是往回拽。方向永远沿绳，而且只拉不推——软绳子不会顶人。一根轻绳各处的拉力处处相等，这是“轻”这个理想化带来的便利。

\textbf{弹力。}弹簧、橡皮筋、蹦床面的力。它是形变最直白的表达：拉得越长（在限度内），回拽得越狠。你在实验第一部分已经量出了它的脾气——伸长多少，力就涨多少，一条直线。这个直来直去的性格马上会被写成一条定律。

\textbf{摩擦力。}两个接触面沿切线方向的互相拖曳。档案要分两页写：

\emph{静摩擦}发生在面与面还没有相对滑动时。它是清单里最特别的一位：\emph{大小按需要供给}，你需要多少它出多少，从零开始，直到一个上限；上限大致和接触面被压紧的程度（也就是支持力）成正比。方向永远对着“将要滑动的趋势”——注意是趋势，不是已经发生的运动。它像什么？它像一位有预算上限的谈判对手：你开价它就奉陪，陪到预算花光就谈崩。它\emph{不像}什么？不像弹簧——弹簧看你压了多深，它不看形变，只看“形势需要”；也不像重力——重力到点上班、雷打不动，静摩擦是你不来它就不出现。

\emph{动摩擦}发生在面与面已经相对滑动时。它的大小近似固定，也大致与支持力成正比，方向对着相对滑动。实验第三部分已经暗示你：动摩擦的价码通常比静摩擦的上限略低，所以东西一旦滑起来就更“收不住”。

\textbf{阻力。}物体穿过空气或水时，流体给的“顶牛”。方向对着相对流体运动的方向，大小随速度长大——慢的时候大约跟速度成正比，快的时候大约跟速度的\emph{平方}成正比，还跟迎着流的面积和流体的密度有关。雨滴就是被它接住的。

这里必须埋一个会让直觉翻车的例子，请先记住它，后面要用：\emph{摩擦力可以和运动同向}。你走路时，脚底向后蹬地，地面给你的静摩擦力是\emph{向前}的——正是它推着你前进；你能在地面上走，全靠这位“谈判对手”肯出钱。卡车启动时货厢里没打滑的箱子，也是被向前的静摩擦力带着加速的。“摩擦总是阻碍运动”这句话错得很经典：摩擦阻碍的是\emph{相对滑动（或滑动趋势）}，而相对滑动的趋势未必和运动反着来。

清单在手，回头看生活，处处是这几位成员在上班：拔河比的不是谁的力气大，而是谁的鞋底和地面之间摩擦预算高——吨位相当的队伍，穿钉鞋的一边几乎稳赢；筷子夹起一粒花生米，靠的是指尖和筷子、筷子和花生之间两份静摩擦的合力；蹦床把你弹回空中，是床面大形变储蓄的弹力在还债；斜拉桥几百吨的桥面悬在半空，是一根根钢索的拉力在轮流值班。学会点名，你就不会再对着一个静止的物体问“它怎么不动”，而会改口问：“此刻都有谁在推它，账是怎么平的？”

清单点完名，还要落到纸上。第 6 章介绍的受力图在这里正式开工：把研究对象单独画出来（划定系统边界），别的什么都不画，只把\emph{别的物体对它}施加的力画成箭头。还有一条纪律要重申：清单上每一位成员都是\emph{相互作用}，都自带第 6 章那对作用力与反作用力——地板顶书柜，书柜也压地板；绳子拉你，你也拽绳子。画受力图时只收“作用在研究对象身上”的那一半，另一半画到对方的图里去，绝不在同一张图里凑对抵消。以那只推而不动的书柜为例：

\begin{center}
\begin{tikzpicture}[>=Stealth,thick]
  \draw[fill=gray!15] (0,0) rectangle (2.2,1.4);
  \node at (1.1,0.7) {书柜};
  \draw (-0.8,0) -- (3.0,0);
  \draw[->,red!70!black] (0,0.7) -- (1.0,0.7) node[midway,above]{$F_{\text{推}}$};
  \draw[->,blue!70!black] (2.2,0.35) -- (1.2,0.35) node[midway,above]{$f_{\text{静}}$};
  \draw[->,green!50!black] (1.1,1.4) -- (1.1,2.4) node[midway,right]{$N$};
  \draw[->,green!50!black] (1.1,0.7) -- (1.1,-0.6) node[midway,right]{$mg$};
\end{tikzpicture}
\end{center}

竖直方向：支持力 \(N\) 与重力 \(mg\) 两清；水平方向：推力 \(F_{\text{推}}\) 与静摩擦 \(f_{\text{静}}\) 两清。合力为零，书柜不动，账平了。注意图里\emph{没有}“向前的运动趋势”这种箭头——受力图只画力，不画运动，更不画“惯性”。画受力图时多画一个子虚乌有的箭头，是全书最常见的画图事故。

\step{一句话模型}

\begin{oneline}
每一种力都来自一种具体的相互作用，各有各的开价方式：重力到点上班，支持力和静摩擦“按需要供给、有上限”，弹力和拉力看形变下菜碟，动摩擦近似固定价，阻力随速度长大。把每种力的脾气摸清，再求合力，剩下的交给第 6 章。
\end{oneline}

\step{数学翻译器}

现在把脾气写成式子。每写一条，立刻用极端情形检查它。

\textbf{胡克定律：弹力正比于形变。}
\[
F=-kx .
\]
它描述什么：弹簧被拉长（或压短）\(x\) 时，往回拽的力有多大。为什么这样写：你的硬币实验画出来的就是一条过原点的直线，直线就长这样；\(k\) 叫劲度系数，是“这根弹簧有多硬”的编号，硬弹簧 \(k\) 大。负号不是大小，是方向：你向右拉，力向左指，永远想把你拽回原长。何时成立：伸长量不太大时。何时失效：实验第二部分——拉过头直线拐弯，橡皮筋永久变长，胡克定律从此不认识这根橡皮筋。

检查：\(x=0\) 时 \(F=0\)，弹簧不拉不压自然不出力，对。\(x\) 加倍 \(F\) 加倍，和你的图一致。方向反转：把“拉”换成“压”，\(x\) 变负，\(F\) 变正，力调头，仍然是“拽回原长”。极小情形：一根 \(k\) 很大的钢棒压短了头发丝直径那么一丁点，力照样可观——支持力的微观真身就在这里。

把实验第一部分的数据画出来，就是这条定律的肖像：

\begin{center}
\begin{tikzpicture}[>=Stealth,scale=0.9]
  \draw[->] (0,0) -- (5.6,0) node[below left]{伸长量 \(x\)};
  \draw[->] (0,0) -- (0,3.4) node[above left]{拉力 \(F\)};
  \draw[thick,blue!70!black] (0,0) -- (3.6,2.7) node[right]{斜率 \(=k\)};
  \draw[dashed,red!70!black] (3.6,2.7) .. controls (4.2,2.85) and (4.6,2.8) .. (5.2,2.4);
  \fill[red!70!black] (3.6,2.7) circle (1.5pt) node[below right]{弹性限度};
  \node[red!70!black] at (4.9,3.1) {拉过头：曲线拐弯};
\end{tikzpicture}
\end{center}

直线段的斜率就是 \(k\)：同一根弹簧，图的陡峭程度不变；换一根更硬的，直线更陡。过了弹性限度，图像拐弯而且不再回头——公式在这一刻作废，不是渐渐不准，而是\emph{换了游戏规则}。

\textbf{静摩擦与动摩擦：一个不等式和一个等式。}
\[
f_{\text{静}}\le \mu_{\text{静}} N,\qquad f_{\text{动}}\approx \mu_{\text{动}} N .
\]
它们描述什么：接触面切向拖曳力的\emph{上限}和滑动时的近似价码。\(\mu\) 叫摩擦因数，由两个面的材料与状态（干燥、上油、结冰）决定，是个要靠实验查的数。为什么静摩擦是不等式：因为它按需供给——左边是实际值，右边是预算上限，实际值由“维持不滑动需要多少”决定，不是由公式直接给出。这是本章最容易写错的地方：把 \(f_{\text{静}}=\mu_{\text{静}}N\) 当等式用，等于假定谈判对手每次都把预算花光。

检查：\(N=0\) 时两者都为零——把书悬空贴着黑板，不压紧就没有摩擦，对。你推书柜的力为零时，静摩擦也是零，不是 \(\mu_{\text{静}}N\)：书柜静静待着，没有理由凭空多出一个横向的力。方向反转：你改为向左推，静摩擦立刻调头向右，它只跟“滑动趋势”对着干，没有自己的偏好。极大情形：推力超过 \(\mu_{\text{静}}N\)，谈判破裂，书柜开滑，价码切换到 \(\mu_{\text{动}}N\)。

摩擦因数要查表才有数：橡胶轮胎对干燥沥青，\(\mu_{\text{静}}\) 大约零点七；木块对木块，大约零点三；钢对冰，不到零点零五。同样是刹车，干地和冰面的制动距离差出十倍——冬天的事故，多半是这条不等式在不同 \(\mu\) 下开出的两张账单。

这个不等式还养活了一项重要的汽车技术：防抱死制动系统（ABS）。车轮\emph{滚动}时，轮胎与地面的接触点相对地面其实是瞬时静止的，起作用的是静摩擦；一旦刹车踩死、车轮抱死打滑，就切换到更小的动摩擦，而且方向盘随之失灵。ABS 每秒十几次地松—紧刹车，让车轮始终待在“将滑未滑”的静摩擦区间——既榨出接近上限的制动力，又保住转向。它不是让刹车更狠，而是让刹车更\emph{懂}摩擦的脾气。

\textbf{阻力：随速度长大的价目表。}
\[
F_{\text{阻}}\approx \tfrac12\,\rho\, C_d\, A\, v^2 \qquad(\text{速度较快时}),
\]
其中 \(\rho\) 是流体密度，\(A\) 是迎风面积，\(C_d\) 是形状系数（流线型小，平板大）。它描述什么：你要以速度 \(v\) 排开前方的空气，每秒得把多少空气推开、推得多猛。为什么有 \(v^2\)：每秒撞上的空气质量正比于 \(v\)，每团空气被你甩开的速度也正比于 \(v\)，两件事相乘，平方就出来了。检查：\(v=0\) 时没有阻力——静止的伞不受风的账，对。速度加倍，阻力变四倍——骑车速二十和四十的区别不是“累一倍”而是“累四倍”，和你的腿的记忆一致。面积加倍阻力加倍——张开伞骑车试试（请在安全的空地上）。汽车工业为这一条公式烧掉了成吨的汽油钱：家用轿车时速一百二时，发动机的大半功率都在对抗空气，所以设计师拼命把车身做成流线、把 \(C_d\) 从零点四压到零点二几，还给车头车尾做风洞实验。省下的每一成 \(C_d\)，都是续航表上实打实的数字。

\textbf{终端速度：账算平的时刻。}雨滴下落时，重力 \(mg\) 向下、阻力向上且随速度长大。速度涨到阻力恰好等于重力时，合力为零，加速度为零，速度从此\emph{不再增长}——不是停住，而是以这个速度匀速下落。令两边相等：
\[
mg=\tfrac12\,\rho\, C_d\, A\, v_t^2
\quad\Longrightarrow\quad
v_t=\sqrt{\dfrac{2mg}{\rho C_d A}} .
\]
检查：质量大的物体 \(v_t\) 大（大的雨滴落得快——第三道猜测题的答案）；迎风面积大 \(v_t\) 小（降落伞的全部原理）；没有空气时 \(\rho=0\)，公式发散——这正是在说“真空中没有终端速度，羽毛和铁锤一起落地”，模型用发散的方式诚实地告诉你它不适用了。

来一个有真实数据的估算。取一颗半径一毫米的雨滴：质量约 \(4\times10^{-6}\) 千克（四毫克），迎风面积 \(A=\pi r^2\approx3\times10^{-6}\) 平方米，\(C_d\) 取 \(0.5\)，空气密度 \(\rho\approx1.2\) 千克每立方米。代入：
\[
v_t=\sqrt{\dfrac{2\times4\times10^{-6}\times9.8}{1.2\times0.5\times3\times10^{-6}}}
\approx\sqrt{44}\approx 7\ \text{米每秒}.
\]
七米每秒，正好就是实测值。一个公式从两千米高空接住了整场雨。

同一公式换个数字就是跳伞运动员：质量七十千克、俯卧迎风面积约零点七平方米，算出 \(v_t\) 约五十五米每秒，接近每小时两百公里——自由落体的跳伞者就是以高速公路的速度在“巡航”；伞一开，面积变成几十平方米，\(v_t\) 跌到五米每秒上下，和从两米高台跳下的冲击相当。降落伞不是“挡住了重力”，而是\emph{把迎风面积做大，让空气早点把账接过去}。

\textbf{圆周运动：转弯是要收钱的。}最后把清单用在转圈上。第 5 章讲过，速度是带方向的量：哪怕速率不变，方向在变，速度就在变，就有加速度。一个物体以速率 \(v\) 沿半径 \(r\) 的圆匀速转圈，它每一瞬的加速度大小是
\[
a=\frac{v^2}{r},
\]
方向永远指向圆心。于是牛顿第二定律立刻要账：必须有大小为 \(mv^2/r\)、指向圆心的\emph{合力}。这笔账就叫“向心力”——但请听清楚：\emph{向心力不是清单上的新成员，它是一个岗位名称}。这个比喻像什么？像医院里的“值班医生”：值班是岗位，内外科的大夫轮流来坐。它不像什么？不像一种新的力种——花名册上查无此人，你不能指着某个物体说“它受到了一个向心力”，只能说“某个（或某几个）真实的力此刻在干向心的活”。谁来上岗都行：汽车转弯，是轮胎与地面的静摩擦上岗；洗衣机甩干，是筒壁对水的支持力上岗（水滴从小孔飞出去，是因为支持力突然撤岗，水沿切线直奔前程，而不是被什么“离心力”甩出去）；月球绕地球，是重力上岗。在受力图上把“向心力”和重力、支持力并列多画一个箭头，等于给同一个岗位发两份工资——这是本章第二大易错点。

\begin{center}
\begin{tikzpicture}[>=Stealth,thick]
  \draw[gray!60] (0,0) circle (1.8);
  \fill (0,0) circle (1.5pt) node[below left=1pt]{圆心};
  \fill[blue!70!black] (1.8,0) circle (2.5pt);
  \draw[->,blue!70!black] (1.8,0) -- (1.8,1.5) node[midway,right]{速度 \(\bm v\)（沿切线）};
  \draw[->,red!70!black] (1.8,0) -- (0.5,0) node[midway,above]{加速度 \(\bm a\)（指向圆心）};
  \fill[blue!70!black] (0,1.8) circle (2.5pt);
  \draw[->,blue!70!black] (0,1.8) -- (-1.3,1.8) node[midway,above]{\(\bm v\)};
  \draw[->,red!70!black] (0,1.8) -- (0,0.6) node[midway,right]{\(\bm a\)};
\end{tikzpicture}
\end{center}

看图记住两件事：速度和加速度\emph{不}同向——一个沿切线，一个指圆心，二者始终垂直；速率可以一点不变，方向却在时时刻刻被“掰弯”，掰弯就要收钱，账是 \(mv^2/r\)。

检查：\(v=0\) 时 \(a=0\)，不转圈不要钱，对。\(r\) 越大 \(a\) 越小：同样的速度，急弯比缓弯凶险得多，开车的直觉完全吻合。\(r\to\infty\) 时 \(a\to0\)：半径无穷大的圆就是直线，匀速直线运动加速度为零，模型平滑地退化成直线情形。

\begin{mathbridge}
\textbf{雨滴是怎么“匀速”起来的。}设阻力在小速度下与速度成正比：\(F_{\text{阻}}=bv\)。取向下为正，牛顿第二定律写成
\[
m\frac{\mathrm dv}{\mathrm dt}=mg-bv .
\]
右边在 \(v=mg/b\) 时等于零，那就是终端速度 \(v_t\)。令 \(u=v-v_t\)，方程变成 \(m\,\mathrm du/\mathrm dt=-bu\)，解是指数衰减：
\[
v(t)=v_t\bigl(1-\mathrm e^{-t/\tau}\bigr),\qquad \tau=\frac{m}{b}.
\]
\(\tau\) 叫特征时间：下落后大约三五个 \(\tau\)，速度就已经很贴近 \(v_t\)。检查：\(t=0\) 时 \(v=0\)，对；\(t\to\infty\) 时 \(v\to v_t\)，对；\(b\to0\) 时 \(\tau\to\infty\)——没有阻力就永远到不了终端，模型再次用发散表达“我不适用”。平方阻力情形思路相同，只是方程不再是线性的，解的曲线形状略不同，结论不变：速度单调逼近一个有限的 \(v_t\)。

\textbf{向心加速度从哪来。}把圆周上相隔很小时间 \(\Delta t\) 的两个速度箭头平移到同一起点，箭头尖端的连线在 \(\Delta t\to0\) 时恰好指向圆心；用弧长与角度的关系 \(v\,\Delta t=r\,\Delta\theta\) 和速度三角形 \(|\Delta\bm v|=v\,\Delta\theta\) 一除，\(a=|\Delta\bm v|/\Delta t=v^2/r\) 就掉出来了。它只用了“速度是向量”这一件事，没有用任何新假设。
\end{mathbridge}

\begin{challenge}
\textbf{摩擦因数为什么是个“凑合数”。}把两个看起来光滑的面放到显微镜下，它们都是山地。真正互相顶住的只是少数山峰，真实接触面积远小于表观面积，而且大致正比于压紧程度——这碰巧解释了 \(f\propto N\) 为什么近似成立。但这个“碰巧”很脆弱：接触面会黏连、会犁出沟、会随速度和温度变化，赛车热熔胎的等效 \(\mu\) 能超过 \(1\)，冰面上薄薄的液态水层会把 \(\mu_{\text{动}}\) 打到很低。所以 \(\mu\) 更像一段经验曲线的斜率，而不是材料的基本常数。真正的微观理论要把接触点的黏附、形变和断裂一个个算清楚，至今仍是活跃的研究领域。

\textbf{“离心力”到底存不存在。}站在地面（近似惯性系）看，旋转的水滴沿切线飞走，全程没有受任何向外的力——“离心力”不存在。但坐进旋转的参考系里看，静止于该系中的物体确实表现出“被向外拉”的倾向，为了在那个参考系里继续使用牛顿定律的形式，数学家给这种倾向起了名字叫\emph{惯性离心力}。它是参考系选择带来的记账项，不是哪个物体发出来的相互作用。两种说法各自自洽，混用才出错——这正是“数学模型”和“物理解释”要分开摆放的又一例。
\end{challenge}

\step{模型的边界}

清单开完了，账也算了，现在给每条公式贴上保质期。

\textbf{胡克定律的边界}你在实验里已经亲手摸到了：弹性限度之外，材料要么永久变形（橡皮筋变长、回形针掰弯），要么直接断。微观上，弹性来自原子间相互作用的“拉开一点就想回去”，那本质上是电磁相互作用——这一笔会在讲电磁和材料的章节里兑现。橡胶其实连线性区都很勉强，它能拉长到几倍，曲线弯得厉害；胡克定律真正的主场是小形变的金属弹簧。

\textbf{摩擦模型的边界}更宽。\(f=\mu N\) 是一个极好用的工程近似，不是自然定律：\(\mu\) 随材料配对、表面污染、湿度、温度、滑动速度而变；静摩擦到动摩擦的切换也不是干净的一刀，而是带着“黏—滑”抖动的过程——琴弓拉出声、粉笔写字、刹车尖叫，都是这种抖动。方向判断永远问“相对滑动趋势朝哪”，而不是“物体朝哪动”。冬天的马路是这个模型失效的展览场：结冰后 \(\mu\) 跳水，于是人们往路面撒盐（降低冰点）、往车轮上缠防滑链（用凸起“咬住”地面，已经不完全是摩擦模型管的事了）。同一个“滑”字，在不同温度和表面下是好几本不同的账。

\textbf{阻力模型的边界}在两头。极慢、极小的物体（比如微尘、油滴）穿过流体时，阻力正比于速度而不是速度平方，那是黏性主导的区域；快到接近声速时，空气开始被剧烈压缩，公式又得换。中间这段“平方律”覆盖了雨滴、跳伞、骑车、汽车巡航，已经够用——这就是为什么它进了本章正文，另外两头的修正留给了流体力学。

\textbf{圆周运动账本的边界}在于“岗位有人肯上”。雨天路面 \(\mu_{\text{静}}\) 变小，\(v^2/r\) 要的账超过摩擦预算，车就向外滑——不是被甩出去，是“要的钱没人出”，车只能沿切线加惯性滑行的曲线走。工程师的对策很妙：把弯道外侧垫高，让路面支持力不再竖直，分一个水平分量出来指向圆心，替摩擦分担账单。你坐高铁过弯几乎不用减速，靠的就是这份“超高”设计。

\textbf{支持力“按需要供给”的边界}是物体得撑得住：椅子有承重上限，桥面有设计载荷。约束不是魔法，是被压紧的材料在用弹性还债，债还不起就塌。

\step{回头解释开场现象}

现在可以拆开雨滴那笔账了。雨滴在云里几乎从静止开始下落，头一两秒确实在加速，越落越快；但阻力跟着速度平方往上蹿，几秒钟之内就追上了重力。从那一刻起，合力为零，雨滴以约七米每秒的终端速度匀速走完剩下的两千米。加速只发生在出发的最初十几米里，全程的百分之九十九都在匀速——所以“两千米高空”和“两百米高空”掉下来的雨，打在脸上几乎没区别。

自由落体公式 \(v=\sqrt{2gh}\) 错了吗？没有。它是一条有保质期的模型：只在“阻力可以忽略”的范围内成立。对一枚从桌上掉落的硬币，阻力占比不到百分之一，公式漂亮；对一颗穿越两千米的雨滴，阻力早已唱主角，直接套用就是拿错了模型——第 1 章说的“模型不是世界本身”，在这里挨了一次最真实的雨淋。

同一本账还顺手解释了一串日常：跳伞者开伞前终端速度约五十五米每秒，开伞后迎风面积暴涨，\(v_t\) 掉到五米每秒上下，刚好是能站稳落地的速度；猫从几层楼跌落常常能走开，而人不行，因为 \(v_t\) 里有质量除以面积，个子越小越占上风；你骑车飙到三十以后举步维艰，因为阻力按平方长、每秒要克服它做的功还要再乘一个速度，身体很快就到极限。一张价目表，管住了从云到猫的所有下落。

最后对一下开头押的三注。第一题答案是（c）：书柜不动，水平方向合力必为零，静摩擦恰好等于你的五十牛——它按需要供给，既不缺席也不超额；选（a）的同学把“没动”当成了“没力”，选（d）的同学还活在“必须更大力才能顶住”的直觉里。第二题答案是（c）：一枚硬币一厘米，三枚大约三厘米，胡克定律的直线就是这么朴实——当然，只在你没把橡皮筋拉过弹性限度时成立。第三题答案是（b）：终端速度 \(v_t=\sqrt{2mg/(\rho C_d A)}\) 里，大雨滴的质量按半径三次方长、面积只按平方长，所以越大落越快——毛毛雨能飘着走，暴雨砸得人生疼，差别全在这个平方根里。第四题答案是（b）：车轮滚动时接地点与地面相对静止，起作用的是预算更高的静摩擦；一旦抱死打滑，就切换到更小的动摩擦，方向还不再听方向盘的。这正是 ABS 干的事，只不过它比任何司机都快。三题全对，说明你的直觉已经完成了一次升级；有猜错的也不亏，错的每一步都是旧模型让位的现场。

\step{脑洞实验与挑战题}

\begin{thought}[把蚂蚁放大一百倍]
一只蚂蚁从摩天楼顶掉下去，拍拍土继续赶路；一匹马从同样高度掉下去，结局惨烈。假设蚂蚁的密度和形状不变，只把尺寸放大一百倍变成“巨蚁”，它从楼顶掉下来会怎样？

算一笔比例账：质量随尺寸的三次方长，放大一百倍，体重变一百万倍；迎风面积只随平方长，变一万倍。终端速度公式 \(v_t=\sqrt{2mg/(\rho C_d A)}\) 里，分子长得比分母快，\(v_t\) 按尺寸的平方根长——巨蚁的终端速度是蚂蚁的十倍。更糟的是落地那一刻：要停住的动量随三次方暴涨，而甲壳和腿的强度（看横截面积）只随平方长。所以童话里的巨型昆虫首先会死在物理学手里，而不是英雄手里。反过来，把大象缩到老鼠大，它从桌边掉下去会毫发无伤。

这笔账反过来也管着天空：为什么鸟类体重的上限大约只有十几千克？翅膀的升力大致随翼面积长，体重却随体积长，越大的鸟越“飞不起自己”，现存最大的飞行动物已经贴着这条红线走。为什么高空跳伞要用伞，雨燕却敢从悬崖直接俯冲起飞？同一平方根，两端各是生死。尺寸的每一次缩放，都在暗地里改写力的账本——这一条会在第 9 章讲能量、在讲材料强度的章节里反复登场。
\end{thought}

\begin{puzzles}
\item 卡车向前加速，货厢里的箱子没有打滑，跟着车一起越来越快。箱子在水平方向只受摩擦力。这个静摩擦指向哪个方向？提示：先问“谁在水平方向上让箱子加速”，再问“如果车厢底板绝对光滑，箱子会怎样”——趋势是判断方向的唯一钥匙。

\item 一根弹簧挂五百克重物，伸长两厘米。现在把弹簧剪成相等的两段，只用其中一段挂同样的重物，伸长量是多少？提示：想象两根“半弹簧”串联：每段只伸长一厘米，合起来才是两厘米。同一根弹簧越短越硬。

\item 雨天汽车过半径固定的弯，工程师给出的安全速度与车的质量几乎无关（假定干燥路面、摩擦因数已知）。为什么质量会约掉？提示：把“需要的账”\(mv^2/r\) 和“摩擦的预算”\(\mu mg\) 并排写，看哪个字母两边都有。

\item 洗衣机脱水时，水珠从衣服的小孔飞出去后，几乎沿直线走。有人说这是“离心力把它沿半径方向甩出去的”。用受力图反驳他。提示：力消失的瞬间，物体保持的是那一刻的\emph{速度}；而那一刻速度正沿圆的切线，不沿半径。
\end{puzzles}

本章的清单到此清点完毕，但清单里第一位的重力，我们只给它立了“地面附近 \(mg\)”这条临时工规定。它凭什么这样？它和让月球绕地球转圈的那个力是什么关系？月亮做圆周运动需要的“向心账”又是谁在付？下一章把这条线追到天上去。
